\subsection{Eichsymmetrie und elektromagnetische Wechselwirkung}
\label{subsec:eichsymmetrie_elektromagnetische_wechselwirkung}

Im vorangegangenen Abschnitt wurde die freie kovariante
Dirac-Gleichung betrachtet. Außerdem wurde das elektromagnetische
Viererpotential als die natürliche relativistische Größe eingeführt,
über die das Dirac-Feld an das elektromagnetische Feld gekoppelt werden
kann.

Noch ist jedoch nicht geklärt, warum diese Kopplung gerade eine
bestimmte mathematische Form besitzen muss. Die zentrale Leitfrage
dieses Abschnitts lautet daher:

\begin{quote}
	\textbf{Warum besitzt die elektromagnetische Wechselwirkung gerade die
	mathematische Form, die in der Quantenelektrodynamik verwendet wird?}
\end{quote}

Die Antwort liegt in einer Symmetrie der Phase des Dirac-Feldes. Die
Forderung nach lokaler \(U(1)\)-Eichsymmetrie führt zunächst zur
Einführung eines \(U(1)\)-Eichfeldes. Identifiziert man die zugehörige
\(U(1)\)-Ladung \(q\) mit der elektrischen Ladung, erhält dieses Eichfeld
die Bedeutung des elektromagnetischen Viererpotentials und beschreibt
die elektromagnetische Kopplung.

\subsubsection*{Eine globale Phasentransformation}

Zunächst wird das Dirac-Feld mit einem überall gleichen komplexen
Phasenfaktor multipliziert:

\begin{equation}
	\psi(x)
	\longrightarrow
	\psi'(x)
	=
	\exp\left(
	\frac{\mathrm{i}q\chi}{\hbar}
	\right)
	\psi(x),
	\qquad
	\chi=\text{konstant}.
	\label{eq:globale_phasentransformation_qed}
\end{equation}

Dabei bezeichnet \(q\) zunächst die zu dieser \(U(1)\)-Transformation
gehörende Ladung. In der QED wird sie mit der elektrischen Ladung
identifiziert; für das Elektronenfeld gilt dann \(q=-e\). Die Größe
\(\chi\) ist hier zunächst eine Konstante.

Eine solche Transformation verändert nur die absolute Phase des
Feldes. Messbare Größen bleiben unverändert, weil sich der Phasenfaktor
mit seinem komplex konjugierten Faktor aufhebt. Für die Dichte gilt
beispielsweise

\begin{equation*}
	\psi'^{\dagger}\psi'
	=
	\exp\left(-\frac{\mathrm{i}q\chi}{\hbar}\right)
	\exp\left(\frac{\mathrm{i}q\chi}{\hbar}\right)
	\psi^{\dagger}\psi
	=
	\psi^{\dagger}\psi.
\end{equation*}

Die absolute Phase einer Wellen- oder Feldfunktion ist daher nicht
messbar. Physikalisch beobachtbar sind nur Größen, in denen sich diese
gemeinsame Phase herauskürzt.

Die Phasenfaktoren

\begin{equation*}
	\exp(\mathrm{i}\alpha),
	\qquad
	\alpha\in\mathbb{R},
	\qquad
	\left|\exp(\mathrm{i}\alpha)\right|=1,
\end{equation*}

bilden die Symmetriegruppe \(U(1)\)\index{U(1)-Symmetrie}. Der Name
bezeichnet hier lediglich die Menge aller komplexen Phasenfaktoren mit
dem Betrag eins. Eine weitergehende Gruppentheorie ist für den
folgenden Gedankengang nicht erforderlich.

\begin{DidacticBox}[Was bedeutet globale \(U(1)\)-Symmetrie?]
	\label{box:globale_u1_symmetrie_qed}

	\emph{Global} bedeutet, dass an jedem Raumzeitpunkt derselbe
	Phasenfaktor verwendet wird.

	Die Transformation verändert die mathematische Darstellung des
	Dirac-Feldes, aber keine physikalische Vorhersage. Die freie
	Dirac-Gleichung besitzt deshalb eine globale
	\(U(1)\)-Symmetrie\index{Eichsymmetrie!globale}.
\end{DidacticBox}

Da \(\chi\) konstant ist, verschwindet ihre Ableitung. Deshalb gilt

\begin{equation}
	\partial_\mu\psi'(x)
	=
	\exp\left(
	\frac{\mathrm{i}q\chi}{\hbar}
	\right)
	\partial_\mu\psi(x).
	\label{eq:ableitung_globale_phase_qed}
\end{equation}

Die freie Dirac-Gleichung erhält dadurch lediglich denselben
Phasenfaktor auf beiden Seiten. Ihre Form und ihre physikalischen
Aussagen bleiben unverändert.

\subsubsection*{Von der globalen zur lokalen Symmetrie}

Nun wird die Symmetrieforderung verschärft. Die Phase soll nicht mehr
überall gleich sein, sondern vom Raumzeitpunkt abhängen:

\begin{equation}
	\chi
	\longrightarrow
	\chi(x).
	\label{eq:phase_wird_lokal_qed}
\end{equation}

Die lokale Phasentransformation\index{Phasentransformation!lokale}
lautet damit

\begin{equation}
	\psi(x)
	\longrightarrow
	\psi'(x)
	=
	\exp\left(
	\frac{\mathrm{i}q\chi(x)}{\hbar}
	\right)
	\psi(x).
	\label{eq:lokale_phasentransformation_qed}
\end{equation}

Jetzt wirkt die Ableitung nicht nur auf \(\psi(x)\), sondern auch auf
den ortsabhängigen Phasenfaktor. Mit der Produkt- und der Kettenregel
folgt

\begin{equation}
	\begin{aligned}
		\partial_\mu\psi'(x)
		&=
		\partial_\mu
		\left[
		\exp\left(
		\frac{\mathrm{i}q\chi(x)}{\hbar}
		\right)
		\psi(x)
		\right]
		\\[4pt]
		&=
		\exp\left(
		\frac{\mathrm{i}q\chi(x)}{\hbar}
		\right)
		\left[
		\partial_\mu\psi(x)
		+
		\frac{\mathrm{i}q}{\hbar}
		\left(\partial_\mu\chi(x)\right)
		\psi(x)
		\right].
	\end{aligned}
	\label{eq:ableitung_lokale_phase_qed}
\end{equation}

Der zweite Term in der eckigen Klammer ist neu. Er verschwindet nur,
wenn \(\chi\) konstant ist. Für eine lokale Phase kann die gewöhnliche
Ableitung daher nicht in derselben einfachen Weise transformieren wie
das Feld selbst.

Setzt man Gleichung~\eqref{eq:ableitung_lokale_phase_qed} in die freie
kovariante Dirac-Gleichung ein, erhält man

\begin{equation}
	\begin{aligned}
	&\left(
	\mathrm{i}\hbar c\,\gamma^\mu\partial_\mu
	-
	m_{\mathrm e}c^2
	\right)
	\psi'(x)
	\\[4pt]
	&\quad=
	\exp\left(
	\frac{\mathrm{i}q\chi(x)}{\hbar}
	\right)
	\left[
	\left(
	\mathrm{i}\hbar c\,\gamma^\mu\partial_\mu
	-
	m_{\mathrm e}c^2
	\right)
	\psi(x)
	-
	qc\,\gamma^\mu
	\left(\partial_\mu\chi(x)\right)
	\psi(x)
	\right].
	\end{aligned}
	\label{eq:freie_diracgleichung_lokal_transformiert_qed}
\end{equation}

Selbst wenn \(\psi(x)\) die freie Dirac-Gleichung erfüllt, bleibt der
zusätzliche letzte Term bestehen. Eine lokale Phasenänderung zerstört
somit zunächst die einfache Form der freien Gleichung.

\subsubsection*{Das Eichfeld und die kovariante Ableitung}

Soll die lokale Phasentransformation dennoch eine Symmetrie sein, muss
der zusätzliche Ableitungsterm ausgeglichen werden. Dazu wird ein
weiteres Feld benötigt: das \(U(1)\)-Eichfeld\index{Eichfeld}
\(A_\mu(x)\).

In Abschnitt~\ref{subsec:freie_wechselwirkende_quantenfelder} wurde das
Viererpotential mit oberem Index als
\(A^\mu=(\phi/c,\mathbf A)\) eingeführt. Mit der im mathematischen
Anhang verwendeten Minkowski-Signatur \((+,-,-,-)\) gilt entsprechend

\begin{equation}
	A_\mu
	=
	\left(
	\frac{\phi}{c},
	-\mathbf A
	\right).
	\label{eq:viererpotential_unterer_index_qed}
\end{equation}

An die Stelle der gewöhnlichen Ableitung tritt nun die
eichkovariante Ableitung\index{kovariante Ableitung!Eichsymmetrie}

\begin{equation}
	D_\mu
	=
	\partial_\mu
	+
	\frac{\mathrm{i}q}{\hbar}A_\mu.
	\label{eq:kovariante_ableitung_qed}
\end{equation}

Diese Definition ist mit den in
Abschnitt~\ref{subsec:freie_wechselwirkende_quantenfelder} verwendeten
SI-Einheiten konsistent. Da \(A^0=\phi/c\) gewählt wurde, besitzen alle
Komponenten von \(A_\mu\) die Einheit eines Impulses pro Ladung. Der Ausdruck
\(qA_\mu/\hbar\) besitzt daher wie \(\partial_\mu\) die Einheit einer
inversen Länge. Ein zusätzlicher Faktor \(c\) ist in
Gleichung~\eqref{eq:kovariante_ableitung_qed} nicht erforderlich.

Der Buchstabe \(D\) erinnert daran, dass nicht mehr nur die räumliche
oder zeitliche Änderung des Feldes gemessen wird. Die Ableitung enthält
zusätzlich den Einfluss des Eichfeldes.

\subsubsection*{Wie sich das Viererpotential transformieren muss}

Damit \(D_\mu\psi\) unter einer lokalen Phasenänderung dieselbe Form
behält wie \(\psi\), muss sich auch das Viererpotential verändern. Die
erforderliche Eichtransformation\index{Eichtransformation!Viererpotential}
lautet

\begin{equation}
	A_\mu(x)
	\longrightarrow
	A'_\mu(x)
	=
	A_\mu(x)
	-
	\partial_\mu\chi(x).
	\label{eq:eichtransformation_viererpotential_qed}
\end{equation}

Mit dieser Transformation lässt sich die gewünschte Eigenschaft direkt
überprüfen:

\begin{equation}
	\begin{aligned}
		D'_\mu\psi'(x)
		&=
		\left(
		\partial_\mu
		+
		\frac{\mathrm{i}q}{\hbar}A'_\mu
		\right)
		\exp\left(
		\frac{\mathrm{i}q\chi(x)}{\hbar}
		\right)
		\psi(x)
		\\[4pt]
		&=
		\exp\left(
		\frac{\mathrm{i}q\chi(x)}{\hbar}
		\right)
		\left[
		\partial_\mu
		+
		\frac{\mathrm{i}q}{\hbar}
		\left(
		A'_\mu
		+
		\partial_\mu\chi
		\right)
		\right]
		\psi(x)
		\\[4pt]
		&=
		\exp\left(
		\frac{\mathrm{i}q\chi(x)}{\hbar}
		\right)
		D_\mu\psi(x).
	\end{aligned}
	\label{eq:transformation_kovariante_ableitung_qed}
\end{equation}

Im letzten Schritt wurde
\(A'_\mu+\partial_\mu\chi=A_\mu\) verwendet. Der störende Zusatzterm
aus Gleichung~\eqref{eq:ableitung_lokale_phase_qed} wird somit genau
durch die Transformation des Eichfeldes aufgehoben.

Für das skalare Potential und das Vektorpotential entspricht
Gleichung~\eqref{eq:eichtransformation_viererpotential_qed} den bekannten
Transformationen

\begin{equation}
	\phi'
	=
	\phi
	-
	\frac{\partial\chi}{\partial t},
	\qquad
	\mathbf A'
	=
	\mathbf A
	+
	\nabla\chi.
	\label{eq:eichtransformation_potentiale_qed}
\end{equation}

Die lokale \(U(1)\)-Eichsymmetrie führt damit zunächst zu einem Eichfeld
mit der Transformationsstruktur eines \(U(1)\)-Viererpotentials.
Identifiziert man die zugehörige Ladung \(q\) mit der elektrischen
Ladung, wird \(A_\mu\) als elektromagnetisches Viererpotential gedeutet
und die Kopplung erhält die Bedeutung der elektromagnetischen
Wechselwirkung.

\subsubsection*{Das Prinzip der minimalen Kopplung}

Der Übergang von der freien zur wechselwirkenden Dirac-Gleichung wird
durch eine einzige Ersetzung vollzogen:

\begin{equation}
	\boxed{
		\partial_\mu
		\longrightarrow
		D_\mu
	}.
	\label{eq:minimale_kopplung_qed}
\end{equation}

Diese Vorschrift heißt minimale Kopplung\index{minimale Kopplung}. Sie
fügt den einfachsten und fundamentalen Wechselwirkungsterm hinzu, der
mit der lokalen \(U(1)\)-Eichsymmetrie vereinbar ist, und bildet die
Grundlage der QED in ihrer üblichen renormierbaren Form.

Wendet man die Ersetzung auf die freie kovariante Dirac-Gleichung aus
Gleichung~\eqref{eq:kovariante_diracgleichung_qed} an, folgt

\begin{equation}
	\left(
	\mathrm{i}\hbar c\,\gamma^\mu D_\mu
	-
	m_{\mathrm e}c^2
	\right)
	\psi(x)
	=
	0.
	\label{eq:diracgleichung_kovariante_ableitung_qed}
\end{equation}

Mit Gleichung~\eqref{eq:kovariante_ableitung_qed} ergibt sich daraus

\begin{equation}
	\left(
	\mathrm{i}\hbar c\,\gamma^\mu\partial_\mu
	-
	qc\,\gamma^\mu A_\mu
	-
	m_{\mathrm e}c^2
	\right)
	\psi(x)
	=
	0.
	\label{eq:diracgleichung_elektromagnetische_kopplung_qed}
\end{equation}

Der neue Term \(-qc\,\gamma^\mu A_\mu\) beschreibt die Kopplung
zwischen dem Dirac-Feld und dem elektromagnetischen Feld. Sein
Vorzeichen enthält die Ladung \(q\) bereits. Für das Elektronenfeld ist
also \(q=-e\) einzusetzen.

In der aus Kapitel VII bekannten Hamiltonform entspricht dieselbe
Ersetzung der Gleichung

\begin{equation}
	\mathrm{i}\hbar
	\frac{\partial\psi}{\partial t}
	=
	\left[
	c\,\boldsymbol{\alpha}\cdot
	\left(
	\hat{\mathbf p}
	-
	q\mathbf A
	\right)
	+
	\beta m_{\mathrm e}c^2
	+
	q\phi
	\right]
	\psi.
	\label{eq:diracgleichung_minimale_kopplung_hamiltonform_qed}
\end{equation}

Damit ist auch in der vertrauten Form sichtbar, dass der Impuls durch
\(\hat{\mathbf p}-q\mathbf A\) und die Energie durch das skalare
Potential \(q\phi\) ergänzt werden.

\begin{PhysicsBox}[Lokale Symmetrie bestimmt die Kopplung]
	\label{box:lokale_symmetrie_bestimmt_kopplung_qed}

	Die Forderung nach lokaler \(U(1)\)-Eichsymmetrie führt zur Einführung
	eines \(U(1)\)-Eichfeldes und der kovarianten Ableitung. Wird die
	zugehörige \(U(1)\)-Ladung \(q\) mit der elektrischen Ladung
	identifiziert, erhält das Eichfeld die Bedeutung des
	elektromagnetischen Viererpotentials.

	\begin{equation*}
		\begin{gathered}
			\text{freie Dirac-Gleichung}
			\\
			+\;\text{lokale }U(1)\text{-Symmetrie}
			\\[4pt]
			\Downarrow
			\\[4pt]
			\partial_\mu\longrightarrow D_\mu
			\quad\text{mit Eichfeld }A_\mu
			\\[4pt]
			\Downarrow
			\\[4pt]
			q=\text{elektrische Ladung}
			\\[4pt]
			\Downarrow
			\\[4pt]
			\text{elektromagnetische Kopplung}.
		\end{gathered}
	\end{equation*}

	Das Symmetrieprinzip bestimmt zunächst die mathematische Struktur des
	\(U(1)\)-Eichfeldes. Durch die Identifikation der \(U(1)\)-Ladung mit
	der elektrischen Ladung wird daraus die elektromagnetische Kopplung.
\end{PhysicsBox}

\subsubsection*{Der elektromagnetische Feldstärketensor}

Das Viererpotential beschreibt das Eichfeld. Die messbaren elektrischen
und magnetischen Felder werden in der relativistischen Schreibweise im
elektromagnetischen Feldstärketensor\index{Feldstärketensor!elektromagnetischer}

\begin{equation}
	F_{\mu\nu}
	=
	\partial_\mu A_\nu
	-
	\partial_\nu A_\mu
	\label{eq:feldstaerketensor_qed}
\end{equation}

zusammengefasst. Die Komponenten mit einem zeitlichen und einem
räumlichen Index enthalten das elektrische Feld. Die rein räumlichen
Komponenten enthalten das magnetische Feld. Eine ausführliche
Komponentendarstellung ist für den weiteren Gedankengang nicht
erforderlich.

Der Feldstärketensor bleibt unter der Eichtransformation des Potentials
unverändert. Einsetzen von
Gleichung~\eqref{eq:eichtransformation_viererpotential_qed} liefert

\begin{equation}
	\begin{aligned}
		F'_{\mu\nu}
		&=
		\partial_\mu
		\left(A_\nu-\partial_\nu\chi\right)
		-
		\partial_\nu
		\left(A_\mu-\partial_\mu\chi\right)
		\\[4pt]
		&=
		F_{\mu\nu}
		-
		\partial_\mu\partial_\nu\chi
		+
		\partial_\nu\partial_\mu\chi
		\\[4pt]
		&=
		F_{\mu\nu}.
	\end{aligned}
	\label{eq:feldstaerketensor_eichinvariant_qed}
\end{equation}

Im letzten Schritt wurde verwendet, dass gewöhnliche partielle
Ableitungen vertauschbar sind. Unterschiedliche Potentiale, die durch
eine Eichtransformation miteinander verbunden sind, beschreiben somit
dieselben elektrischen und magnetischen Felder.

\subsubsection*{Von der Symmetrie zur vollständigen Theorie}

Die Argumentation dieses Abschnitts lässt sich als geschlossene Kette
zusammenfassen:

\begin{equation}
	\begin{aligned}
		\text{globale Phasensymmetrie}
		&\longrightarrow
		\text{lokale Phasensymmetrie}
		\\
		&\longrightarrow
		\text{zusätzlicher Ableitungsterm}
		\\
		&\longrightarrow
		\text{Eichfeld }A_\mu
		\longrightarrow
		D_\mu
		\\
		&\longrightarrow
		\text{Identifikation }q=\text{elektrische Ladung}
		\\
		&\longrightarrow
		\text{elektromagnetische Wechselwirkung}.
	\end{aligned}
	\label{eq:leitstruktur_eichsymmetrie_qed}
\end{equation}

Damit ist erklärt, wie die lokale \(U(1)\)-Eichsymmetrie zu einem
\(U(1)\)-Eichfeld und zur kovarianten Ableitung führt. Identifiziert man
die \(U(1)\)-Ladung mit der elektrischen Ladung, beschreibt diese
Struktur die elektromagnetische Wechselwirkung; die minimale Kopplung
liefert dabei den fundamentalen Wechselwirkungsterm der üblichen
renormierbaren QED.

Für die vollständige Quantenelektrodynamik fehlt jedoch noch eine
gemeinsame mathematische Struktur, die das freie Dirac-Feld, das freie
elektromagnetische Feld und ihre Wechselwirkung zugleich enthält.
Genau diese Aufgabe übernimmt die Lagrangedichte der
Quantenelektrodynamik, die im nächsten Abschnitt eingeführt wird.
